\chapter{Dodatkowe przykłady i uwagi}
\label{cha:dodatkowe_przyklady}

\section{Liczby i jednostki}
Liczby i jednostki zapisujemy czcionką tekstową, niezależnie czy występują w tekście, czy w równaniu.

Do zapisu liczb i jednostek używamy poleceń \texttt{\textbackslash num}, \texttt{\textbackslash unit} oraz \texttt{\textbackslash qty} z pakietu \texttt{siunitx}. Część dziesiętną w tym wypadku zawsze zapisujemy po kropce. W zależności od ustawień języka, kropka pozostanie kropką lub zostanie zamieniona na przecinek. Pakiet zadba też o poprawne separatory tysięczne, znaki minus oraz odstęp pomiędzy liczbą a jednostką.

\begin{lstlisting}
    Liczba $\pi$ w przybliżeniu wynosi \num{3.14}.
    Jednostką częstotliwości jest \unit{\hertz}.
    Stała Boltzmanna w przybliżeniu wynosi \qty{1.38e-23}{\joule\per\kelvin}.
    Impedancja wyjściowa to \complexqty{10-j30}{\ohm}.
    Powierzchnia układu wynosi \qtyproduct{10 x 20}{\micro\meter}.
    Pasmo FR3 obejmie częstotliwości pomiędzy FR1 i FR2, tj.
        \qtyrange{7.125}{24.25}{\giga\hertz}.
\end{lstlisting}

    Liczba $\pi$ w przybliżeniu wynosi \num{3.14}.

    Jednostką częstotliwości jest \unit{\hertz}.

    Stała Boltzmanna w przybliżeniu wynosi \qty{1.38e-23}{\joule\per\kelvin}.

    Impedancja wyjściowa to \complexqty{10-j30}{\ohm}.

    Powierzchnia układu wynosi \qtyproduct{10 x 20}{\micro\meter}.

    Pasmo FR3 obejmie częstotliwości pomiędzy FR1 i FR2, tj. \qtyrange{7.125}{24.25}{\giga\hertz}.

    \vskip1em

    Więcej przykładów w \textit{User manual} na stronie pakietu \url{https://ctan.org/pkg/siunitx}.

\section{Symbole}

W przeciwieństwie do liczb i jednostek, symbole fizyczne, matematyczne (np. zmienne), elektryczne (np. oznaczenia elementów na schemacie), itp. zapisujemy czcionką matematyczną. Czcionką tekstową zapisujemy natomiast nazwy funkcji matematycznych (pełne i skrócone), skrótowce oraz słowa (pełne i skrócone).

Jeśli symbol występuje w równaniu, środowisko domyślnie zmieni jego krój. Jeśli symbol umieszczamy w tekście, należy otoczyć go znakami dolara.

Jeśli symbol posiada indeks (górny lub dolny), to krój indeksu również powinien być dobrany według powyższych reguł. Wyjątkowo, jeśli indeks jest jednoliterowy, nawet jeśli jest skrótem, może być pochylony, o ile nie prowadzi to niekonsekwencji w zapisie podobnych symboli. Zazwyczaj zatem, krój indeksu powinien być:

\begin{description}[labelwidth=2.5cm, leftmargin=2.5cm]
    \item[prosty] jeśli jest liczbą (np. $f_\text{0}$), słowem ($t_\text{stop}$), akronimem lub skrótem (np. $V_\text{in}$); wówczas obudowujemy go środowiskiem \texttt{\textbackslash text\{\}}
    \item[pochyły] jeśli jest symbolem (np. $\Delta_t$, $g_m$).
\end{description}

\textbf{Przykład:}

Tranzystor $M_\text{1}$ wymiarach $W$/$L$\,$=$\,\qty{20}{\micro\meter}/\qty{40}{\nano\meter} znajduje się w zakresie nasycenia ($V_{DS} > V_{GS} - V_{TH}$), wówczas prąd płynący przez jego dren można wyrazić następująco:
\begin{equation}
    I_{D\text{1}} = \frac{\num{1}}{\num{2}}\mu_{n}C_\text{ox}\frac{W}{L}\left(V_{GS}-V_{TH}\right)^2
\end{equation}

\section{Inne uwagi}
    Polski cudzysłów w \LaTeX u otwieramy dwoma przecinkami \texttt{{,}{,}} a zamykamy dwoma apostrofami \texttt{{'}{'}} --- w ten sposób otrzymamy ,,tekst w cudzysłowie''.

    Każda nowa myśl powinna stanowić nowy akapit. Warto stosować płynne przejścia tak, żeby zakończenie jednego akapitu naturalnie prowadziło do rozpoczęcia następnego.


%----------------------------------------------------------------------
